% Chapter 3 --- Integrable cones and the Pettis integral
% (paper §4, milestone M3).
%
% Mirrors theories/icones/:
%   pettis.v             paper Def 4.1 — Pettis-style cone integral
%   icone.v              paper §4 — integrable-cone HB tower
%   icone_integral.v     paper Lemma 4.2, 4.6, 4.7, Thm 4.5, 4.12, 4.15
%   fubini.v             paper Thm 4.15 — cone-Fubini, paper-form Fubini
%   icone_cat.v          paper Thm 4.16, 4.18 — ICones complete + well-powered
%   examples_icone.v     canonical examples

%%%%%%%%%%%%%%%%%%%%%%%%%%%%%%%%%%%%%%%%%%%%%%%%%%%%%%%%%%%%%%%%%%%%%%%%
\chapter{Integrable cones}\label{ch:icones}

This chapter formalises paper \S4: the Pettis-style integral, the
$\ICone$ \HB{} tower, cone-Fubini, and the categorical structure of
$\ICones$ (complete and well-powered). All material in this chapter
lives under \texttt{theories/icones/}.

The novelty of paper \S4 is that the cone-valued integral of a
measurable path $\beta : X \to B$ against a finite measure
$\mu \in \FMeas(X)$ is \emph{not} built up from simple-function
approximations, as in the Lebesgue case. Instead, following
Pettis~\cite{Pet38}, it is characterised by a universal property:
$\Pettis{X}{\beta\,d\mu}$ is the unique element $x \in B$ such that
every arity-$0$ test $m \in \Mtest_B$ satisfies
$m(x) = \int_X m(\beta(r))\,\mu(dr)$. Uniqueness is immediate from
the separation axiom (Mssep) on $B$; existence is the defining
condition of an integrable cone (paper Def~4.3). The formalisation
mirrors this Pettis-by-axiom structure: the integral value is
extracted via classical indefinite description from the existence
witness packaged as the $\isICone$ \HB{} mixin.

\section{The Pettis-style integral}\label{sec:pettis}

\begin{definition}[Pettis integral specification]\label{def:path-integral-eq}
\uses{def:fmeas, def:path}
Let $B$ be a measurable cone, $X \in \Ar$, $\beta : X \to B$ a path
and $\mu \in \FMeas(X)$. An element $x \in B$ is said to satisfy the
\emph{Pettis equation} for $(\beta, \mu)$ when, for every arity-$0$
test $m \in \Mtest_B$ at the singleton arity $\Ar_0$,
\[
  m(x) \;=\; \int_X m(\beta(r))\,\mu(dr).
\]
The right-hand side is the ordinary Lebesgue integral of the
non-negative bounded measurable function $m \circ \beta$ against the
finite measure $\mu$, viewed in the formalisation via the
$\fine$-cast of its $\overline{\R}$-valued integral. In Rocq this is
the predicate \rocq{Icones.icones.pettis.path_integral_eq}.
\rocqok
\end{definition}

\begin{lemma}[Uniqueness of the Pettis integral]\label{lem:path-integral-eq-unique}
\uses{def:path-integral-eq}
Two elements $x_1, x_2 \in B$ that both satisfy the Pettis equation
for $(\beta, \mu)$ are equal:
\[
  m(x_1) = m(x_2) \;\text{for every } m \in \Mtest_B \;\Longrightarrow\; x_1 = x_2,
\]
by the separation axiom (Mssep) on $B$. In Rocq this is
\rocq{Icones.icones.pettis.path_integral_eq_unique}. \rocqok
\end{lemma}

\begin{proof}
Both equations evaluate the same Lebesgue integral, so the two test
values $m(x_1)$ and $m(x_2)$ agree for every $m \in \Mtest_B$ at
arity $0$. The (Mssep) axiom of $\MCone$ then forces $x_1 = x_2$.
\end{proof}

By Lemma~\ref{lem:path-integral-eq-unique}, when an element $x$
satisfying the Pettis equation exists, it is unique. We denote it by
$\Pettis{X}{\beta\,d\mu}$, mirroring the paper's
$\int_X \beta(r)\mu(dr)$. In Rocq, once the existence witness $\hex$
is in hand, the value is extracted via classical indefinite
description as $\rocq{Icones.icones.pettis.path_integral}\ \hex$.

\section{Integrable cones}\label{sec:icones}

\begin{definition}[Integrable cone --- paper Def 4.3]\label{def:icone}
\uses{def:path-integral-eq}
A measurable cone $B$ is \emph{integrable} when, for every
$X \in \Ar$, every measurable path $\beta : X \to B$ and every
finite measure $\mu \in \FMeas(X)$, there exists $x \in B$
satisfying the Pettis equation for $(\beta, \mu)$.

In the formalisation this is the \HB{} mixin
\rocq{Icones.icones.icone.isICone} layered over $\MCone$. The
existence field is named \texttt{icone\_exists}; the corresponding
structure \rocq{Icones.icones.icone.ICone} packages an $\MCone$ with
this Pettis-witness axiom. We write $\iconeType\ \Ar$ for the type of
integrable cones over a given $\Ar$-parameter. \rocqok
\end{definition}

For an integrable cone $B$, every triple $(X, \beta, \mu)$ with
$\beta$ a measurable path determines a unique element of $B$
satisfying the Pettis equation; we package it as the Rocq function
\rocq{Icones.icones.icone.icone_integral}, with specification
\rocq{Icones.icones.icone.icone_integralP} and uniqueness
\rocq{Icones.icones.icone.icone_integral_eqP}. We write
$\Pettis{X}{\beta\,d\mu}$ for $\rocq{Icones.icones.icone.icone_integral}\
\beta\ H_\beta\ \mu$.

\section{The norm bound}\label{sec:integral-norm-bound}

\begin{lemma}[Paper Lemma 4.2]\label{lem:path-integral-norm-le}
\uses{def:path-integral-eq}
Let $\beta : X \to B$ be a measurable path bounded in cone-norm by
$M_\beta \geq 0$, and let $\mu \in \FMeas(X)$. Any element $x \in B$
satisfying the Pettis equation for $(\beta, \mu)$ obeys
\[
  \cnorm{x}_B \;\leq\; M_\beta \cdot \cnorm{\mu}_{\FMeas(X)}.
\]
In Rocq this is \rocq{Icones.icones.icone_integral.path_integral_norm_le}.
\rocqok
\end{lemma}

\begin{proof}
If $x = 0$ the bound is trivial. Otherwise, for any $\varepsilon > 0$
the axiom (Msnorm) on $B$ provides $m \in \Mtest_B \setminus \{0\}$
at arity $0$ with $\cnorm{x} \leq \varepsilon + m(x)/\cnorm{m}$.
Combining the Pettis equation with the pointwise test bound
$m(\beta(r)) \leq \cnorm{m} \cdot \cnorm{\beta(r)} \leq
\cnorm{m} \cdot M_\beta$ and integrating against $\mu$ yields
$m(x) = \int_X m(\beta(r))\,\mu(dr) \leq \cnorm{m} \cdot M_\beta
\cdot \cnorm{\mu}$. Dividing by $\cnorm{m}$ and letting
$\varepsilon \to 0$ gives the announced inequality.
\end{proof}

\section{The kernel-integral measurability lemma}\label{sec:lemma46}

The next lemma is the technical prerequisite for proving that paths
of paths integrate to paths. It is the only place in M3 where the
mathcomp-analysis library on finite kernels is invoked directly.

\begin{lemma}[Paper Lemma 4.6]\label{lem:kernel-integral-measurable}
\uses{def:fmeas}
Let $\varphi : Y \times X \to \R_{\geq 0}$ be measurable and bounded
(say $\varphi \leq M_\varphi$), and let $\kappa : Y \to \FMeas(X)$
be a bounded kernel (measurable set-evaluations, uniformly bounded
total mass). Then the partial-integral function
\[
  s \in Y \;\longmapsto\; \int_X \varphi(s, r)\,\kappa(s)(dr) \;\in\;
  \R_{\geq 0}
\]
is measurable. In Rocq this is
\rocq{Icones.icones.icone_integral.kernel_integral_measurable}
(real-valued, \fine{}-cast form) and
\rocq{Icones.icones.icone_integral.kernel_integral_measurable_ereal}
(extended-real form). \rocqok
\end{lemma}

\begin{proof}
The paper proves this by monotone-class reduction: the property
holds for simple $\varphi$ and lifts to general $\varphi$ via
monotone convergence (each non-negative measurable $\varphi$ is the
pointwise sup of an increasing sequence of simple functions). The
formalisation routes through the mathcomp-analysis HB factory
\texttt{Kernel\_isFinite}: the bounded kernel $\kappa$ is packaged
as an $R$-finite kernel, and the
\texttt{measurable\_fun\_integral\_finite\_kernel} lemma supplies
the measurability of $s \mapsto \int \varphi(s, r)\,\kappa(s)(dr)$
in the $\overline{\R}$-valued form. The $\fine$-cast to $\R$ is
then measurable because the integrand is bounded by $M_\varphi$ and
$\kappa$ has uniformly finite total mass.
\end{proof}

\section{Lemma 4.7 --- the workhorse}\label{sec:lemma47}

Paper Lemma~4.7 is the engine of all downstream constructions: the
integration operator $I^B_X(\beta, \mu) := \Pettis{X}{\beta\,d\mu}$
is bilinear, $\omega$-continuous and \emph{jointly measurable} on
$\Path(X, B) \mathbin{\&} \FMeas(X) \to B$. We split it into four
sub-statements matching the structure of the Rocq sources.

\subsection{Bilinearity in $\beta$ and $\mu$}\label{ssec:bilinear}

\begin{lemma}[Linearity in $\beta$]\label{lem:integral-lin-beta}
\uses{def:icone}
For every $\mu \in \FMeas(X)$, the integration operator
$\beta \mapsto \Pettis{X}{\beta\,d\mu}$ is $\R_{\geq 0}$-linear:
\begin{align*}
  \Pettis{X}{(\beta_1 + \beta_2)\,d\mu} &= \Pettis{X}{\beta_1\,d\mu}
    + \Pettis{X}{\beta_2\,d\mu}, \\
  \Pettis{X}{(r \cdot \beta)\,d\mu} &= r \cdot \Pettis{X}{\beta\,d\mu}
    \quad (r \in \R_{\geq 0}).
\end{align*}
In Rocq, the predicate-level statements are
\rocq{Icones.icones.icone_integral.path_integral_eq_addB} and
\rocq{Icones.icones.icone_integral.path_integral_eq_scaleB}; the
$\iconeType$ wrap-ups are
\rocq{Icones.icones.icone_integral.icone_integral_addB} and
\rocq{Icones.icones.icone_integral.icone_integral_scaleB}. \rocqok
\end{lemma}

\begin{proof}
Apply the linearity field \texttt{test\_linD} (resp.\
\texttt{test\_linZ}) of $\Mtest_B$ to the Pettis equations of
$\beta_1$, $\beta_2$ (resp.\ $\beta$). The cone-side identity for
the candidate witness $x_1 + x_2$ (resp.\ $r \cdot x$) then matches
the linear combination of the Lebesgue integrals by
$\mathsf{ge0\_integralD}$ (resp.\ $\mathsf{ge0\_integralZl}$) of
mathcomp-analysis, both of which require the finiteness of the
integrand --- which follows from the path bound and (Mssep) is
\emph{not} invoked here. The $\iconeType$ wrap-up applies
$\mathsf{icone\_integral\_eqP}$ (uniqueness) to identify the
candidate with the integral.
\end{proof}

\begin{lemma}[Linearity in $\mu$]\label{lem:integral-lin-mu}
\uses{def:icone}
For every measurable path $\beta : X \to B$, the integration
operator $\mu \mapsto \Pettis{X}{\beta\,d\mu}$ is
$\R_{\geq 0}$-linear:
\begin{align*}
  \Pettis{X}{\beta\,d(\mu_1 + \mu_2)} &= \Pettis{X}{\beta\,d\mu_1}
    + \Pettis{X}{\beta\,d\mu_2}, \\
  \Pettis{X}{\beta\,d(r \cdot \mu)} &= r \cdot \Pettis{X}{\beta\,d\mu}
    \quad (r \in \R_{\geq 0}).
\end{align*}
In Rocq, the predicate-level statements are
\rocq{Icones.icones.icone_integral.path_integral_eq_addmu} and
\rocq{Icones.icones.icone_integral.path_integral_eq_scalemu}; the
$\iconeType$ wrap-ups are
\rocq{Icones.icones.icone_integral.icone_integral_addmu} and
\rocq{Icones.icones.icone_integral.icone_integral_scalemu}. \rocqok
\end{lemma}

\begin{proof}
The measure-side analogue of the previous lemma. The
$\FMeas(X)$-side $+$ is the measure-sum
$\mathsf{measure\_add}$ on the underlying $\overline{\R}$-valued
measures, and $r \cdot \mu$ is mathcomp-analysis's $\mathsf{mscale}$.
The key identities are
$\mathsf{ge0\_integral\_measure\_add}$ and
$\mathsf{ge0\_integral\_mscale}$, applied to the bounded
measurable integrand $r \mapsto m(\beta(r))$.
\end{proof}

\subsection{$\omega$-continuity in $\beta$ and $\mu$}\label{ssec:omega-cont}

\begin{lemma}[$\omega$-continuity in $\beta$]\label{lem:integral-omega-cont-path}
\uses{def:icone, lem:path-integral-norm-le}
Let $(\beta_n)_{n \in \N}$ be an increasing chain of measurable
paths from $X$ to $B$, all bounded by $1$ in cone-norm, with
pointwise supremum $\beta_\infty := \sup_n \beta_n$ a measurable
path. Let $\mu \in \FMeas(X)$ with $\cnorm{\mu} \leq 1$. Then the
chain $\bigl(\Pettis{X}{\beta_n\,d\mu}\bigr)_n$ is increasing and
bounded, with unit-ball supremum
\[
  \bigsup_n \Pettis{X}{\beta_n\,d\mu}
  \;=\; \Pettis{X}{\beta_\infty\,d\mu}.
\]
In Rocq this is
\rocq{Icones.icones.icone_integral.integral_omega_cont_path}.
\rocqok
\end{lemma}

\begin{proof}
The chain $(\Pettis{X}{\beta_n\,d\mu})_n$ is increasing by the
integral-chain-monotonicity helper
\rocq{Icones.icones.icone_integral.icone_integral_chain_le}, which
itself uses the precone residue extraction
\rocq{Icones.icones.icone_integral.precone_residue} to package the
pointwise difference $\beta_{n+1} - \beta_n$ as a measurable path and
appeal to additivity (Lemma~\ref{lem:integral-lin-beta}). The norm
bound on each integral comes from
Lemma~\ref{lem:path-integral-norm-le} with $M_\beta = 1$. The two
sides agree by (Mssep): testing against any arity-$0$ $m$, both
sides equal $\int_X m(\beta_\infty(r))\,\mu(dr)$, by the
mathcomp-analysis monotone convergence theorem
$\mathsf{monotone\_convergence}$ applied to the non-decreasing
sequence $r \mapsto m(\beta_n(r))$.
\end{proof}

\begin{lemma}[$\omega$-continuity in $\mu$]\label{lem:integral-omega-cont-meas}
\uses{def:icone, lem:path-integral-norm-le}
Let $\beta : X \to B$ be a measurable path bounded by $1$, and let
$(\mu_n)_n$ be an increasing chain of finite measures with
$\cnorm{\mu_n} \leq 1$. Then
$\bigl(\Pettis{X}{\beta\,d\mu_n}\bigr)_n$ is increasing and bounded,
with unit-ball supremum
\[
  \bigsup_n \Pettis{X}{\beta\,d\mu_n}
  \;=\; \Pettis{X}{\beta\,d(\bigsup_n \mu_n)}.
\]
In Rocq this is
\rocq{Icones.icones.icone_integral.integral_omega_cont_meas}.
\rocqok
\end{lemma}

\begin{proof}
By the measure-side additivity Lemma~\ref{lem:integral-lin-mu} the
integral chain is non-decreasing, and the norm bound is again
Lemma~\ref{lem:path-integral-norm-le}. To identify the sup, test
against any arity-$0$ $m$ and apply the project-internal
measure-direction monotone convergence theorem
\rocq{Icones.icones.icone_integral.integral_meas_sup}: this version
of MCT (in the \emph{measure} argument rather than the function
argument) is not present in mathcomp-analysis 1.16, and is proved
from $\mathsf{ge0\_integral\_measure\_series}$ via the telescoping
$\mathsf{fmeas\_dseq}$ decomposition of the chain (cf.\ \S3.2.1).
(Mssep) on $B$ then identifies the cone-side sup with the
integral against the limit measure.
\end{proof}

\subsection{Joint measurability}\label{ssec:joint-meas}

\begin{lemma}[Joint measurability]\label{lem:integral-joint-meas}
\uses{def:icone, lem:kernel-integral-measurable}
Let $(S, \Sigma_S)$ be a measurable space, let $\beta : S \to X \to B$
be a family of measurable paths with $(s, r) \mapsto m(z, \beta(s)(r))$
jointly measurable on $\Ar(Z) \times S \times X$ for every test
$m \in \Mtest^Z_B$, and let $\kappa : S \to \FMeas(X)$ be a bounded
kernel (each set-evaluation $s \mapsto \kappa(s)(U)$ measurable, with
uniform mass bound). Then for every $Z \in \Ar$ and every test
$m \in \Mtest^Z_B$, the function
\[
  (z, s) \in \Ar(Z) \times S \;\longmapsto\;
  m\bigl(z, \Pettis{X}{\beta(s)\,d\kappa(s)}\bigr)
\]
is jointly measurable. In Rocq this is
\rocq{Icones.icones.icone_integral.icone_integral_joint_measurable}.
\rocqok
\end{lemma}

\begin{proof}
A general arity-$Z$ test reduces to the arity-$0$ Pettis equation
via the constant-reindexing helper
\rocq{Icones.icones.icone_integral.icone_integral_test_pettis}: at
each $z$, $m(z, \cdot)$ acts like an arity-$0$ test on $B$ obtained
by precomposition with the constant $\Ar_0 \to Z$ at $z$. Hence the
test value at the integral equals the partial integral
$\fine\bigl(\int_X (m(z, \beta(s)(r)))^\sharp\, \kappa(s)(dr)\bigr)$,
which is jointly measurable in $(z, s)$ by
Lemma~\ref{lem:kernel-integral-measurable} (taking $Y := \Ar(Z) \times S$
and $\varphi((z, s), r) := m(z, \beta(s)(r))$).
\end{proof}

This is the form used by every downstream consumer: it powers the
$\iconeType$ instance on $\Path(X, B)$ (\S\ref{sec:path-icone}) and
the joint test-measurability needed by the product and equaliser
structures in $\ICones$ (\S\ref{sec:icones-cat}).

\section{$\FMeas(X)$ is integrable --- paper Thm 4.5}\label{sec:fmeas-icone}

\begin{theorem}[Paper Thm 4.5]\label{thm:fmeas-icone}
\uses{def:icone, def:fmeas}
For every $X \in \Ar$, the measurable cone $\FMeas(X)$ is integrable.
\rocqok
\end{theorem}

\begin{proof}
Let $Y \in \Ar$, $\kappa \in \Path(Y, \FMeas(X))$ a bounded kernel,
and $\nu \in \FMeas(Y)$. The candidate integral is the
\emph{pushforward of paths}: the set function
\[
  \mu : U \in \sigma_X \;\longmapsto\; \int_Y \kappa(s)(U)\,\nu(ds)
  \;\in\; \R_{\geq 0}.
\]
Standard measure theory shows $\mu$ is a finite measure on $X$, and
by the very definition of $\Mtest_{\FMeas(X)}$ in terms of the
generator tests $e_U : \rho \mapsto \rho(U)$, the integral $\mu$
satisfies the Pettis equation for $(\kappa, \nu)$. In Rocq the
existence witness is built in \texttt{examples\_icone.v} as
\rocq{Icones.icones.examples_icone.fmeas_int_exists}; the supporting
work is in \texttt{fmeas\_int\_fun}, \texttt{fmeas\_int\_finP},
\texttt{fmeas\_int\_canon} and \texttt{fmeas\_int\_pettis}.
\end{proof}

\section{$\Path(X, B)$ is integrable --- paper Thm 4.12}%
\label{sec:path-icone}

\begin{theorem}[Paper Thm 4.12]\label{thm:path-icone}
\uses{def:icone, lem:integral-joint-meas}
If $B$ is an integrable cone, then for every $X \in \Ar$ the
measurable cone $\Path(X, B)$ is integrable. The integral is
\emph{pointwise}: if $\eta : Y \to \Path(X, B)$ is a measurable path
of paths and $\nu \in \FMeas(Y)$, the integral $\beta \in \Path(X, B)$
of $\eta$ over $\nu$ is given by
\[
  \beta(r) \;=\; \Pettis{Y}{\eta(s)(r)\,d\nu(s)} \quad (r \in X).
\]
In Rocq this is captured by the \HB.instance at
\rocq{Icones.icones.examples_icone.path_int_exists}. \rocqok
\end{theorem}

\begin{proof}
The candidate $\beta(r) := \Pettis{Y}{\eta(s)(r)\,d\nu(s)}$ exists
pointwise because $B$ is integrable. Its cone-norm bound
$\cnorm{\beta(r)} \leq \cnorm{\eta}\,\cnorm{\nu}$ follows from
Lemma~\ref{lem:path-integral-norm-le}. Measurability of $\beta$ as
a path in $\Path(X, B)$ reduces to a joint test-measurability
statement, which is precisely
Lemma~\ref{lem:integral-joint-meas} applied with $S := X$ and the
swapped roles $(\eta, \nu, r) \rightsquigarrow (\beta(\cdot)(r), \nu, \cdot)$.
The Pettis equation for $\beta$ is then a direct unfolding of the
$\Path$ tests $p = r \rhd m$, since
$p(\beta) = m(\beta(r)) = \int_Y m(\eta(s)(r))\,\nu(ds) = \int_Y p(\eta(s))\,\nu(ds)$.
The full \HB.instance landed via M3 wave~7, after the $\ar\text{-}\mathsf{prod}$
cast-measurability machinery (\S\ref{sec:cone-fubini}) made the joint
measurability hypothesis dischargeable from a plain
$\mathsf{is\_measurable\_path}\ \eta$.
\end{proof}

\section{The cone-Fubini theorem --- paper Thm 4.15}%
\label{sec:cone-fubini}

\begin{theorem}[Paper Thm 4.15, cast-free form]\label{thm:fubini-cone-eq}
\uses{def:icone, lem:integral-joint-meas}
Let $B$ be an integrable cone, $X, Y \in \Ar$, $\mu \in \FMeas(X)$,
$\nu \in \FMeas(Y)$, and let $\beta : X \times Y \to B$ be uniformly
norm-bounded, sectionally measurable in each variable, and jointly
test-measurable in the two iteration shapes
$(z, x, y) \mapsto m(z, \beta(x, y))$ and
$(z, y, x) \mapsto m(z, \beta(x, y))$ for every test
$m \in \Mtest^Z_B$. Then the two iterated cone-integrals agree:
\[
  \Pettis{X}{\Bigl(\Pettis{Y}{\beta(x, y)\,d\nu(y)}\Bigr)\,d\mu(x)}
  \;=\;
  \Pettis{Y}{\Bigl(\Pettis{X}{\beta(x, y)\,d\mu(x)}\Bigr)\,d\nu(y)}.
\]
In Rocq this is \rocq{Icones.icones.fubini.fubini_cone_eq}. \rocqok
\end{theorem}

\begin{proof}
By (Mssep) on $B$, it suffices to check the equation under every
arity-$0$ test $m \in \Mtest_B$. Testing each side, the Pettis
equation applied twice yields
$\fine\!\int_X \fine\!\int_Y m(\beta(x, y))\,d\nu\,d\mu$ on one
side and
$\fine\!\int_Y \fine\!\int_X m(\beta(x, y))\,d\mu\,d\nu$ on the
other. The non-negative bounded measurable function
$m \circ \beta : X \times Y \to \R$ then falls under the standard
mathcomp-analysis Fubini--Tonelli theorems
$\mathsf{fubini\_tonelli1}$ and $\mathsf{fubini\_tonelli2}$, applied
to the product measure $\fmeas\text{-}\mathsf{fin\_view}\,\mu \times
\fmeas\text{-}\mathsf{fin\_view}\,\nu$. The finiteness of every
inner integral (used by the $\fine$ cast) is supplied by the
uniform bound $M_\beta$ together with $\cnorm{\mu}, \cnorm{\nu} < \infty$.
\end{proof}

\begin{theorem}[Paper Thm 4.15, $\ar\text{-}\mathsf{prod}$ form]%
\label{thm:fubini-cone-eq-arprod}
\uses{thm:fubini-cone-eq}
Let $\beta : \ar\text{-}\mathsf{prod}(\Ar, X, Y) \to B$ be a
measurable path. Pulling back along the (measurable) cast
$\ar\text{-}\mathsf{prod\text{-}cast}$ gives a $X \times Y$-indexed
$\beta_{\!\mathsf{uncast}} = \beta \circ \ar\text{-}\mathsf{prod\text{-}cast}$
to which Theorem~\ref{thm:fubini-cone-eq} applies, after deriving
its five hypotheses from $\mathsf{is\_measurable\_path}\,\beta$. In
Rocq this is the paper-form wrapper
\rocq{Icones.icones.fubini.fubini_cone_eq_arprod}. \rocqok
\end{theorem}

\begin{proof}
The cast-measurability field $\ar\text{-}\mathsf{prod\text{-}cast\text{-}meas}$
of $\MeasSubcat$ (introduced in paper \S3, wave 7 of M3 closed it
cleanly) makes $\beta_{\!\mathsf{uncast}}$'s norm bound, section
measurability and joint test-measurability all derivable; the
helpers
\rocq{Icones.icones.fubini.β_uncast_bound},
\rocq{Icones.icones.fubini.β_uncast_secY},
\rocq{Icones.icones.fubini.β_uncast_secX},
\rocq{Icones.icones.fubini.β_uncast_jointX} and
\rocq{Icones.icones.fubini.β_uncast_jointY} discharge each side
condition.
\end{proof}

\section{$\ICones$ is complete and well-powered}\label{sec:icones-cat}

The category $\ICones$ has integrable cones as objects, and as
morphisms those $\MCones$-morphisms $f : B \to C$ that
\emph{preserve integrals}: for every $X, \beta, \mu$,
$f(\Pettis{X}{\beta\,d\mu}) = \Pettis{X}{f(\beta)\,d\mu}$. The
Rocq packaging is \rocq{Icones.icones.icone_cat.icones_hom}; the
category structure (identity, composition, associativity, units) is
in the section \texttt{IConesCat}.

\begin{theorem}[Paper Thm 4.16]\label{thm:icones-complete}
\uses{def:icone, lem:integral-joint-meas}
The category $\ICones$ is complete: it has all small products and
all binary equalisers (hence all small limits, since these generate
all limits in any category).
\end{theorem}

\begin{proof}
\emph{Small products.} Let $(C_i)_{i \in I}$ be a family of
integrable cones with product $P = \mathbin{\&}_{i \in I} C_i$ in
$\Cones$. The test family on $P$ is generated by the lifted tests
$\mathsf{iniTest}_i(m)(r, \vec{x}) := m(r, x_i)$ for
$m \in \Mtest^Y_{C_i}$; the closure axioms (Mscomp), (Mssep) and
(Msnorm) carry over from each $C_i$ (uniform existence-of-test for
(Msnorm) uses that $\cnorm{\vec{x}} = \sup_i \cnorm{x_i}$).
Integrability is componentwise: for $\gamma \in \Path(X, P)$ and
$\mu \in \FMeas(X)$, the candidate is
$\vec{x} := (\Pettis{X}{(\gamma_i(r))\,d\mu(r)})_{i \in I}$, in
$P$ by the norm bound (Lemma~\ref{lem:path-integral-norm-le}), and
the Pettis equation holds by definition of the test family.
Projections and tuplings preserve paths and integrals; the universal
property is exactly the $\Cones$-tupling restricted to integrable
cones. In Rocq the construction lives in section
\texttt{IConesProducts} of \texttt{icone\_cat.v}; the integrable cone
structure is the HB.instance at
\texttt{icones\_prod\_int\_exists}, packaged as
\rocq{Icones.icones.icone_cat.icones_prod}, with projection
\rocq{Icones.icones.icone_cat.icones_proj} and tupling
\rocq{Icones.icones.icone_cat.icones_tuple} (universal property
\rocq{Icones.icones.icone_cat.icones_tuple_proj},
\rocq{Icones.icones.icone_cat.icones_tuple_unique}).

\emph{Binary equalisers.} Let $f, g \in \ICones(C, D)$ and let
$(P, e \in \Cones(P, C))$ be the equaliser of $f, g$ in $\Cones$. The
test family on $E := P$ inherits from $C$ by precomposition with the
inclusion $e$, satisfying (Msnorm) since $B_P = P \cap B_C \subseteq
B_C$ tightens the test sup. Integrability: for $\beta \in \Path(X, E)$
and $\mu \in \FMeas(X)$, the $C$-integral
$x := \Pettis{X}{e(\beta(r))\,d\mu(r)}$ satisfies $f(x) = g(x)$
(both equal $\Pettis{X}{f(\beta(r))\,d\mu(r)} =
\Pettis{X}{g(\beta(r))\,d\mu(r)}$ since $f, g$ preserve integrals
and $\beta$ lands in $E$), hence $x \in E$ and is the $E$-integral
of $\beta$. In Rocq, section \texttt{IConesEqualisers},
\HB.instance at \texttt{icones\_eq\_int\_exists}, the equaliser
object \rocq{Icones.icones.icone_cat.icones_eq}, its inclusion
\rocq{Icones.icones.icone_cat.icones_eq_incl} and the mediating
arrow \rocq{Icones.icones.icone_cat.icones_eq_med} (universal
property \rocq{Icones.icones.icone_cat.icones_eq_med_factor},
\rocq{Icones.icones.icone_cat.icones_eq_med_unique}).
\end{proof}

\begin{theorem}[Paper Thm 4.18]\label{thm:icones-well-powered}
\uses{def:icone}
In the category $\ICones$:
\begin{enumerate}
  \item $1 = \R_{\geq 0}$ is both a separator and a coseparator;
  \item $\ICones$ is well-powered.
\end{enumerate}
\end{theorem}

\begin{proof}
\emph{Coseparator.} For $f \neq g$ in $\ICones(B, C)$ and $x \in B$
with $f(x) \neq g(x)$, (Mssep) on $C$ yields an arity-$0$ test
$m \in \Mtest^0_C$ with $m(f(x)) \neq m(g(x))$. Each such $m$ is
itself a morphism $C \to 1$ in $\ICones$ (it is linear,
$\omega$-continuous, non-expansive, and the Pettis equation
\emph{is} the integral-preservation equation against $m$). In Rocq:
\rocq{Icones.icones.icone_cat.icones_coseparator}.

\emph{Separator.} For each $x \in B$, the map $b_x(\lambda) := \lambda
\cdot x$ from $\R_{\geq 0} \cong 1$ to $B$ is in $\ICones(1, B)$:
linearity and $\omega$-continuity are immediate, and the
integral-preservation equation $b_x(\Pettis{X}{\beta\,d\mu}) =
\Pettis{X}{b_x(\beta)\,d\mu}$ comes from the linearity of the
ordinary Lebesgue integral (test against any arity-$0$
$m \in \Mtest_B$ and use $m(b_x(\lambda)) = \lambda \cdot m(x)$).
Since $b_x(1) = x$, the family $\{b_x\}_{x \in B}$ separates parallel
morphisms. In Rocq:
\rocq{Icones.icones.icone_cat.icones_separator}.

\emph{Well-poweredness.} Let $h \in \ICones(D, C)$ be a monomorphism.
Since $1$ is a separator, $h$ is injective on points
(\rocq{Icones.icones.icone_cat.icones_subobject_inj}). By the
isomorphism lemma~4.17 of the paper (every bijection from an
integrable cone transports an integrable-cone structure), one may
replace $D$ by its image $S := h(D) \subseteq C$ equipped with the
transported structure, so that $h$ is identified with the set-inclusion
$S \hookrightarrow C$. The class of structures on a fixed subset
$S \subseteq C$ is bounded by the cardinality of the structure
space
$S^{S \times S} \times S^{\R_{\geq 0} \times S} \times \R_{\geq 0}^S
\times \prod_{X \in \Ar} \mathcal{P}\bigl(\R_{\geq 0}^{X \times S}\bigr)$,
which is a set because $\Ar$ is small. Hence the class of
subobjects of $C$ (up to $\ICones$-isomorphism) injects into the
small type
$\{(S, F) \mid S \subseteq C\}$, witnessed in Rocq by the bound
$\rocq{Icones.icones.icone_cat.icones_well_powered_bound} := \mathsf{set}\ C$.
\end{proof}

\section{Design choices}\label{sec:design-choices}

The formalisation makes a handful of deliberate choices that
diverge from the paper's prose but are equivalent in content.

\paragraph{Pettis-by-axiom.}
Rather than building $\Pettis{X}{\beta\,d\mu}$ constructively (which
would require equipping every $\iconeType$ with a $\mathsf{choiceType}$
structure on top of its $\Cones$ data), we extract the integral via
classical indefinite description (Rocq's $\mathsf{cid}$) from the
existence witness packaged in the
$\rocq{Icones.icones.icone.isICone}$ mixin. Uniqueness via
(Mssep) (\rocq{Icones.icones.pettis.path_integral_eq_unique}) then
makes the extracted value canonical up to propositional equality.
This keeps the integrable-cone tower lightweight: no choice
structure is required on the cone carriers, only on the ambient
classical metatheory.

\paragraph{Measure-side monotone convergence.}
mathcomp-analysis 1.16 ships $\mathsf{monotone\_convergence}$ in
the \emph{function} argument but not in the \emph{measure}
argument. The $\omega$-continuity-in-$\mu$ proof
(Lemma~\ref{lem:integral-omega-cont-meas}) needs the latter, so we
prove it in-tree as
\rocq{Icones.icones.icone_integral.integral_meas_sup}, via the
telescoping decomposition $\mathsf{fmeas\_dseq}$ of an increasing
chain into a series (cf.\ paper \S3.2.1) plus
$\mathsf{ge0\_integral\_measure\_series}$.

\paragraph{Precone residue.}
The chain-monotonicity helper for integrals
(\rocq{Icones.icones.icone_integral.icone_integral_chain_le})
needs to extract a measurable-path residue $r \mapsto z_r$ such that
$\beta_2(r) = \beta_1(r) + z_r$ when $\beta_1 \cle \beta_2$
pointwise. We package the pointwise extraction in
\rocq{Icones.icones.icone_integral.precone_residue}: a thin wrapper
over $\mathsf{cid}$ that turns the $\PCone$-order predicate
$\cle$ into its sigma-type companion $\{z \mid y = x + z\}$.
Measurability of the residue follows from $\mathsf{test\_linD}$ and
$\mathsf{measurable\_funB}$.

\paragraph{Cast-free Fubini.}
Paper Theorem~4.15 is naturally stated as an equation between
integrals over $X \times Y$ on one side and iterated integrals on
the other. In our formalisation $\Ar(X \times Y) =
\ar\text{-}\mathsf{prod}(\Ar, X, Y)$ has a non-trivial
measurability cast to $\Ar(X) \times \Ar(Y)$. The natural
cast-free formulation
\rocq{Icones.icones.fubini.fubini_cone_eq} is the equation of the
two iterated integrals, which is what every downstream consumer
needs. The paper form
\rocq{Icones.icones.fubini.fubini_cone_eq_arprod} is recovered as a
thin wrapper once the cast-measurability field
$\ar\text{-}\mathsf{prod\text{-}cast\text{-}meas}$ of $\MeasSubcat$
is in place (M3 wave~7).
